\documentclass[11pt,a4paper]{article}

% =========================================================
% PACKAGES
% =========================================================

\usepackage[
    a4paper,
    left=2cm,
    right=2cm,
    top=1.6cm,
    bottom=1.6cm
]{geometry}

\usepackage[T1]{fontenc}
\usepackage{lmodern}
\usepackage{graphicx}
\usepackage{xcolor}
\usepackage{array}
\usepackage{tabularx}
\usepackage{ragged2e}
\usepackage{enumitem}
\usepackage{fancyhdr}
\usepackage{hyperref}

% =========================================================
% GENERAL SETTINGS
% =========================================================

\setlength{\parindent}{0pt}
\setlength{\parskip}{0pt}

\definecolor{headingblue}{RGB}{60,95,140}

% =========================================================
% HEADER AND FOOTER
% =========================================================

\pagestyle{fancy}

\fancyhf{}

% BTECH and Team ID kept close together
\fancyhead[L]{
    \textbf{\textit{BTECH}}
    \hspace{1.5cm}
    \textbf{DSCPP-III-2026-T201}
}

\fancyfoot[R]{
    \textit{Page \thepage}
}

\renewcommand{\headrulewidth}{0pt}
\renewcommand{\footrulewidth}{0pt}

\setlength{\headheight}{15pt}
\setlength{\headsep}{20pt}
\setlength{\footskip}{22pt}

% =========================================================
% HEADING COMMANDS
% =========================================================

\newcommand{\mainheading}[1]{%
    \vspace{2mm}
    {\color{headingblue}
    \bfseries
    \fontsize{16}{19}\selectfont
    #1}
    \vspace{5mm}
}

\newcommand{\sectionheading}[1]{%
    {\color{headingblue}
    \fontsize{14}{17}\selectfont
    #1}
    \vspace{2mm}
}

% =========================================================
% TEXT BOX
% =========================================================

\newcommand{\textanswer}[1]{%
    \vspace{2mm}
    \noindent
    \fbox{%
        \begin{minipage}{0.965\textwidth}
            \vspace{2mm}
            \justifying
            \fontsize{10.5}{13}\selectfont
            #1
            \vspace{2mm}
        \end{minipage}
    }
    \vspace{3mm}
}

% =========================================================
% DOCUMENT
% =========================================================

\begin{document}

% =========================================================
% PAGE 1
% =========================================================

\mainheading{PROJECT AND TEAM INFORMATION}

% ---------------------------------------------------------
% PROJECT TITLE
% ---------------------------------------------------------

\sectionheading{Project Title}

\textanswer{
\textbf{Smart Music Player and Playlist Management System}
}

\vspace{7mm}

% ---------------------------------------------------------
% STUDENT / TEAM INFORMATION
% ---------------------------------------------------------

\sectionheading{Student / Team Information}

\vspace{2mm}

% =========================================================
% TEAM TABLE
% =========================================================

\renewcommand{\arraystretch}{1.1}

\noindent
\begin{tabularx}{\textwidth}{|p{0.62\textwidth}|X|}
\hline

\textbf{Team Name:}\\
\textit{Team \#}
&
\textit{SoundSync}
\\
\hline

% ---------------------------------------------------------
% TEAM MEMBER 1
% ---------------------------------------------------------

\textbf{Team member 1 (Team Lead)}\\
\textit{(Last Name, name: student ID: email, picture)}
&
\begin{minipage}[t]{\linewidth}
\vspace{2mm}

\textit{Goel,Sanya}\\
\textit{2510010795}\\
\textit{2510010795@geu.ac.in}

\vspace{2mm}

\begin{center}
\includegraphics[
    width=2.7cm,
    height=3.0cm,
    keepaspectratio
]{S.png}
\end{center}

\vspace{1mm}
\end{minipage}
\\
\hline

% ---------------------------------------------------------
% TEAM MEMBER 2
% ---------------------------------------------------------

\textbf{Team member 2}\\
\textit{(Last Name, name: student ID: email, picture)}
&
\begin{minipage}[t]{\linewidth}
\vspace{2mm}

\textit{Arora,Manvi}\\
\textit{2510010232}\\
\textit{2510010232@geu.ac.in}

\vspace{2mm}

\begin{center}
\includegraphics[
    width=2.7cm,
    height=3.0cm,
    keepaspectratio
]{ME IMG.jpeg}
\end{center}

\vspace{1mm}
\end{minipage}
\\
\hline

% ---------------------------------------------------------
% TEAM MEMBER 3
% ---------------------------------------------------------

\textbf{Team member 3}\\
\textit{(Last Name, name: student ID: email, picture)}
&
\begin{minipage}[t]{\linewidth}
\vspace{2mm}

\textit{Ram,Eshu}\\
\textit{2510011472}\\
\textit{2510011472@geu.ac.in}

\vspace{2mm}

\begin{center}
\includegraphics[
    width=2.7cm,
    height=3.0cm,
    keepaspectratio
]{ESHU IMG.png}
\end{center}

\vspace{1mm}
\end{minipage}
\\
\hline

\end{tabularx}

% =========================================================
% PAGE 2
% =========================================================

\newpage

\mainheading{PROPOSAL DESCRIPTION }

% ---------------------------------------------------------
% MOTIVATION
% ---------------------------------------------------------

\sectionheading{Motivation }

\textanswer{
Music players are commonly used for listening to songs, creating playlists,
and managing recently played and upcoming songs. Our project aims to develop
a simple C++-based music player containing a predefined collection of 20--25
songs. Users can enter a song name, search the available collection, view the
song name, artist name, album name, and album artwork, and play the selected
audio. The project will provide a practical application of Data Structures
and Object-Oriented Programming concepts.
}

% ---------------------------------------------------------
% STATE OF THE ART
% ---------------------------------------------------------

\sectionheading{State of the Art / Current solution }

\textanswer{
Currently, music applications provide features such as song searching,
playlists, playback, and listening history. Commercial platforms contain
very large music libraries and advanced features. Our proposed system will
implement a smaller music player using a fixed collection of songs. It will
focus on applying C++ programming, OOP concepts, Data Structures, searching,
sorting, and file handling to create a simple and functional music player.
}

% ---------------------------------------------------------
% PROJECT GOALS
% ---------------------------------------------------------

\sectionheading{Project Goals and Milestones }

\textanswer{
The main goal is to develop a simple C++ music player that allows users to
search and play songs from a predefined collection. The system will display
the song name, artist, album and album artwork. C++ concepts such as classes,
objects, encapsulation, inheritance, abstraction and polymorphism will be used
to organize the system.

The proposed milestones are: finalize requirements and song collection,
design the data structures and C++ class structure, implement the music
library and searching, implement playlist and playback features, integrate
the UI and audio playback, and test the complete system.
}

% ---------------------------------------------------------
% PROJECT APPROACH
% ---------------------------------------------------------

\sectionheading{Project Approach }

\textanswer{
The project will be developed in C++, with the main focus on applying Data
Structures and Algorithms to a music-player system. The system will be
divided into modules for maintaining the song library, searching, playlists,
playback, and history. OOP concepts such as classes, objects, encapsulation,
abstraction, inheritance, and polymorphism will be used to structure these
modules.

The proposed data structures include an Array/Vector for storing the 20--25
songs, Linked List for playlists, Stack for recently played songs, Queue for
upcoming songs, and Hash Table for fast song searching.
}

% =========================================================
% PAGE 3
% SYSTEM ARCHITECTURE
% =========================================================

\newpage

\sectionheading{System Architecture (High Level Diagram)}

\vspace{7mm}

\begin{center}

\includegraphics[
    width=0.72\textwidth,
    height=19.5cm,
    keepaspectratio
]{new.jpeg}

\end{center}

% =========================================================
% PAGE 4
% =========================================================

\newpage


% ---------------------------------------------------------
% PROJECT OUTCOME
% ---------------------------------------------------------

\sectionheading{Project Outcome / Deliverables }

\textanswer{
The expected outcome is a working C++-based prototype of a Smart Music Player
containing 20--25 predefined songs. The system will allow users to search for
songs, display their song name, artist, album and album artwork, and play the
selected audio.

The implementation will demonstrate C++ Object-Oriented Programming through
classes and objects representing songs, playlists and the music player. It
will also demonstrate the practical application of arrays/vectors, linked
lists, stacks, queues, hashing, trees, graphs, searching and sorting.

The main deliverables will include the C++ source code, implemented data
structures, test cases, system architecture diagram, project report,
presentation, and final demonstration.
}

% ---------------------------------------------------------
% ASSUMPTIONS
% ---------------------------------------------------------

\sectionheading{Assumptions}

\textanswer{
The system will use a predefined collection of 20--25 legally obtained audio
files and corresponding album images. Users are assumed to search for songs
available in the defined collection. The initial system will be a C++-based
prototype with a simple user interface and local audio playback. The project
will focus primarily on demonstrating Data Structures, Algorithms, and
Object-Oriented Programming concepts rather than providing the complete
functionality of a commercial music platform.
}

% ---------------------------------------------------------
% REFERENCES
% ---------------------------------------------------------

\sectionheading{References}

\vspace{2mm}

\noindent
\fbox{
\begin{minipage}{0.965\textwidth}
\vspace{2mm}

\begin{itemize}[
    leftmargin=6mm,
    itemsep=2pt,
    topsep=0pt
]

\item \url{https://claude.com/}

\item \url{https://www.canva.com/}

\item \url{https://www.google.com/}

\item \url{https://chatgpt.com/}

\item \url{https://www.overleaf.com/}

\item \url{https://www.latex-project.org/}

\end{itemize}

\vspace{1mm}
\end{minipage}
}

% =========================================================
% END
% =========================================================

\end{document}